
\de{ĐỀ THI HỌC KỲ II NĂM HỌC 2022-2023}{THPT Phan Ngọc Hiển - Cà Mau}
\begin{center}
	\textbf{PHẦN 1 - TRẮC NGHIỆM}
\end{center}
\Opensolutionfile{ans}[ans/ans]
%Câu 1
\begin{ex}%[0T9Y4-5] %[Dự án đề kiểm tra HKII NH22-23- Phạm Văn Long]%[Phan Ngọc Hiển - Cà Mau]
	Phương trình nào sau đây là phương trình chính tắc của một hypebol?
	\choice
	{$\dfrac{x^2}{16} - \dfrac{y^2}{9} = -1$}
	{$\dfrac{x^2}{25} + \dfrac{y^2}{9} = 1$}
	{\True $\dfrac{x^2}{9} - \dfrac{y^2}{16} = 1$}
	{$\dfrac{x^2}{16} + \dfrac{y^2}{25} = 1$}
	\loigiai{
		Phương trình $\dfrac{x^2}{9} - \dfrac{y^2}{16} =1 $ là phương trình chính tắc của một hypebol.
	}
\end{ex}

%Câu 2
\begin{ex}%[0T8B3-1]%[Dự án đề kiểm tra HKII NH22-23- Phạm Văn Long]%[Phan Ngọc Hiển - Cà Mau]
	Viết khai triển theo công thức nhị thức Newton của $\left(x^2 - y\right)^5$.
	\choice
	{$x^{10} + 5x^8y + 10x^6y^2 + 10x^4y^3 + 5x^2y^4 + y^5$}
	{\True $x^{10} - 5x^8y + 10x^6y^2 - 10x^4y^3 + 5x^2y^4 - y^5$}
	{$x^{10} + 5x^8y - 10x^6y^2 + 10x^4y^3 - 5x^2y^4 + y^5$}
	{$x^{10} - 5x^8y - 10x^6y^2 - 10x^4y^3 - 5x^2y^4 + y^5$}
	\loigiai{
		Theo khai triển của nhị thức Newton, ta có
		\allowdisplaybreaks
		\begin{eqnarray*}
			\left(x^2 - y\right)^5 
			&=& \mathrm{C}^0_5 \left(x^2\right)^5 + \mathrm{C}^1_5 \left(x^2\right)^4\left(-y\right) + \mathrm{C}^2_5 \left(x^2\right)^3\left(-y\right)^2 + \mathrm{C}^3_5\left(x^2\right)^2\left(-y\right)^3 + \mathrm{C}^4_5 x^2(-y)^4 + \mathrm{C}^5_5 (-y)^5\\
			&=& x^{10} - 5x^8y + 10x^6y^2 - 10x^4y^3 + 5x^2y^4 - y^5.\\
		\end{eqnarray*}
	}
\end{ex}

%Câu 3
\begin{ex}%[0T8B1-2]%[Dự án đề kiểm tra HKII NH22-23- Phạm Văn Long]%[Phan Ngọc Hiển - Cà Mau]
	Từ các số $1$, $2$, $3$, $4$, $5$, $6$ lập được bao nhiêu số tự nhiên gồm $4$ chữ số?
	\choice
	{$2019$}
	{$1297$}
	{\True $1296$}
	{$2110$}
	\loigiai{
		Số số tự nhiên có $4$ chữ số lập được từ các số $1$, $2$, $3$, $4$, $5$, $6$ là $6^4 = 1296$ số.
	}
\end{ex}

%Câu 4
\begin{ex}%[0T9B2-1]%[Dự án đề kiểm tra HKII NH22-23- Phạm Văn Long]%[Phan Ngọc Hiển - Cà Mau]
	Phương trình tham số của đường thẳng $(d)$ đi qua $M(-2;3)$ và có véc-tơ chỉ phương $\overrightarrow{u} = (1;-4)$ là
	\choice
	{$\heva{&x=1 - 2t\\&y=-4 + 3t}$}
	{\True $\heva{&x=-2 +t \\&y=3-4t}$}
	{$\heva{&x=3- 2t \\&y=-4+t}$}
	{$\heva{&x=-2+3t \\&y=1+4t}$}
	\loigiai{
		$(d)$ có véc-tơ chỉ phương $\overrightarrow{u}=(1;-4)$ và đi qua $M(-2;3)$ nên có phương trình tham số là $\heva{&x = -2 +t\\&y=3 -4t.}$
	}
\end{ex}

%Câu 5
\begin{ex}%[0T0B2-2]%[Dự án đề kiểm tra HKII NH22-23- Phạm Văn Long]%[Phan Ngọc Hiển - Cà Mau]
	Gieo ngẫu nhiên một con súc xắc. Xác suất để lần đầu gieo xuất hiện mặt $6$ chấm là
	\choice
	{$\dfrac{1}{3}$}
	{\True $\dfrac{1}{6}$}
	{$\dfrac{1}{2}$}
	{$\dfrac{5}{6}$}
	\loigiai{
		Xác suất để lần đầu gieo xuất hiện mặt $6$ chấm là $\mathrm{P}=\dfrac{1}{6}$.
	}
\end{ex}

%Câu 6
\begin{ex}%[0T9B2-3]%[Dự án đề kiểm tra HKII NH22-23- Phạm Văn Long]%[Phan Ngọc Hiển - Cà Mau]
	Xác định vị trí tương đối của hai đường thẳng $\Delta_1\colon x-2y+1=0$ và $\Delta_2\colon -3x + 6y -10=0$.
	\choice
	{Trùng nhau}
	{Cắt nhau nhưng không vuông góc}
	{Vuông góc nhau}
	{\True Song song}
	\loigiai{
		Vì $\dfrac{1}{-3}=\dfrac{-2}{6}\neq \dfrac{1}{-10}$ nên $(\Delta_1)\parallel (\Delta_2)$.
	}
\end{ex}

%Câu 7
\begin{ex}%[0T0B2-3]%[Dự án đề kiểm tra HKII NH22-23- Phạm Văn Long]%[Phan Ngọc Hiển - Cà Mau]
	Một hộp đựng $4$ bi xanh và $6$ bi đỏ. Chọn ngẫu nhiên $2$ viên bi. Xác suất để rút được $1$ bi xanh và $1$ bi đỏ là
	\choice
	{$\dfrac{4}{15}$}
	{$\dfrac{6}{25}$}
	{$\dfrac{2}{15}$}
	{\True $\dfrac{8}{15}$}
	\loigiai{
		Xác suất để rút được $1$ bi xanh và $1$ bi đỏ là $\mathrm{P}=\dfrac{\mathrm{C}^1_4\cdot \mathrm{C}^1_6}{\mathrm{C}^2_{10}}=\dfrac{8}{15}$.
	}
\end{ex}

%Câu 8
\begin{ex}%[0T0B2-3]%[Dự án đề kiểm tra HKII NH22-23- Phạm Văn Long]%[Phan Ngọc Hiển - Cà Mau]
	Một bình đựng $5$ quả cầu xanh, $4$ quả cầu đỏ và $3$ quả cầu vàng. Chọn ngẫu nhiên $3$ quả cầu. Xác suất để được $3$ quả cầu khác màu là
	\choice
	{$\dfrac{3}{5}$}
	{$\dfrac{3}{14}$}
	{$\dfrac{3}{7}$}
	{\True $\dfrac{3}{11}$}
	\loigiai{
		Xác suất để được $3$ quả cầu khác màu là $\mathrm{P}=\dfrac{\mathrm{C}^1_5\cdot \mathrm{C}^1_4 \cdot \mathrm{C}^1_3}{\mathrm{C}^3_{12}}=\dfrac{3}{11}$.
	}
\end{ex}

%Câu 9
\begin{ex}%[0T8Y2-2]%[Dự án đề kiểm tra HKII NH22-23- Phạm Văn Long]%[Phan Ngọc Hiển - Cà Mau]
	Công thức tính số chỉnh hợp chập $k$ của $n$ phần tử là
	\choice
	{$\mathrm{A}^k_n = \dfrac{n!}{(n-k)!k!}$}
	{$\mathrm{C}^k_n = \dfrac{n!}{(n-k)!k!}$}
	{$\mathrm{C}^k_n = \dfrac{n!}{(n-k)!}$}
	{\True $\mathrm{A}^k_n = \dfrac{n!}{(n-k)!}$}
	\loigiai{
		Công thức tính số chỉnh hợp chập $k$ của $n$ phần tử là $\mathrm{A}^k_n = \dfrac{n!}{(n-k)!}$.
	}
\end{ex}

%Câu 10
\begin{ex}%[0T8K1-2]%[Dự án đề kiểm tra HKII NH22-23- Phạm Văn Long]%[Phan Ngọc Hiển - Cà Mau]
	Từ các chữ số $1$; $2$; $3$; $4$; $5$; $6$. Có thể lập được bao nhiêu số tự nhiên chẵn có $3$ chữ số khác nhau?
	\choice
	{$3!$}
	{\True $60$}
	{$120$}
	{$32$}
	\loigiai{
		Gọi số có $3$ chữ số có dạng $\overline{abc}$ ($a$, $b$, $c \in\left\{1;2;3;4;5;6\right\}$).
		\begin{itemize}
			\item Chọn $c$ có $3$ cách $(c\in \{2;4;6\})$.
			\item Chọn $a$ có $5$ cách.
			\item Chọn $b$ có $4$ cách.
		\end{itemize}
		Vậy có $3\cdot 5\cdot 4=60$ số tự nhiên chẵn có $3$ chữ số khác nhau.
	}
\end{ex}

%Câu 11
\begin{ex}%[0T9B2-5]%[Dự án đề kiểm tra HKII NH22-23- Phạm Văn Long]%[Phan Ngọc Hiển - Cà Mau]
	Tính khoảng cách $\mathrm{d}$ từ điểm $M(1;-3)$ đến đường thẳng $\Delta\colon 3x-4y +10=0$.
	\choice
	{$\mathrm{d}=3$}
	{$\mathrm{d}=2$}
	{$\mathrm{d}=10$}
	{\True $\mathrm{d}=5$}
	\loigiai{
		Khoảng cách $\mathrm{d}$ từ $M(1;-3)$ đến đường thẳng $\Delta\colon 3x-4y +10=0$ là
		\[\mathrm{d}\left(M,\Delta\right)=\dfrac{\left|3\cdot 1 -4\cdot(-3) + 10\right|}{\sqrt{3^2 + (-4)^2}}=\dfrac{25}{5}=5.\]
	}
\end{ex}

%Câu 12
\begin{ex}%[0T9B4-1]%[Dự án đề kiểm tra HKII NH22-23- Phạm Văn Long]%[Phan Ngọc Hiển - Cà Mau]
	Trong mặt phẳng $(Oxy)$, cho elip $(E)$ có phương trình $\dfrac{x^2}{36} + \dfrac{y^2}{9}=1$. Một tiêu điểm của $(E)$ là
	\choice
	{\True $\left(-\sqrt{27};0\right)$}
	{$\left(0;27\right)$}
	{$\left(27;0\right)$}
	{$\left(0;\sqrt{27}\right)$}
	\loigiai{
		Với elip $(E) \colon \dfrac{x^2}{36} + \dfrac{y^2}{9}=1$, ta có $a^2=36$, $b^2=9$.\\
		Suy ra $c = \sqrt{a^2 - b^2} = \sqrt{36-9}=\sqrt{27}$.\\
		Vậy $(E)$ có hai tiêu điểm là $F_1(-\sqrt{27};0)$ và $F_2(\sqrt{27};0)$.
	}
\end{ex}

%Câu 13
\begin{ex}%[0T8B1-2]%[Dự án đề kiểm tra HKII NH22-23- Phạm Văn Long]%[Phan Ngọc Hiển - Cà Mau]
	Từ thành phố $A$ đến thành phố $B$ có $4$ con đường, từ thành phố $B$ đến thành phố $C$ có $3$ con đường, từ thành phố $C$ đến thành phố $D$ có $2$ con đường. Hỏi có bao nhiêu cách đi từ $A$ đến $D$ mà qua $B$ và $C$ chỉ một lần?
	\choice
	{$18$}
	{$9$}
	{$10$}
	{\True $24$}
	\loigiai{
		Số cách đi từ $A$ đến $D$ mà qua $B$ và $C$ chỉ một lần là $4\cdot 3\cdot 2 = 24$ cách.
	}
\end{ex}

%Câu 14
\begin{ex}%[0T0B1-1]%[Dự án đề kiểm tra HKII NH22-23- Phạm Văn Long]%[Phan Ngọc Hiển - Cà Mau]
	Gieo con súc sắc hai lần. Biến cố $A$ là biến cố để sau hai lần gieo có ít nhất một mặt $6$ chấm.
	\choice
	{\True $A=\left\{(1,6);(2,6);(3,6);(4,6);(5,6);(6,6);(6,1);(6,2);(6,3);(6,4);(6,5)\right\}$}
	{$A=\left\{(1,6);(2,6);(3,6);(4,6);(5,6);(6,6)\right\}$}
	{$A=\left\{(6,1);(6,2);(6,3);(6,4);(6,5)\right\}$}
	{$A=\left\{(1,6);(2;6);(3;6);(4;6);(5;6)\right\}$}
	\loigiai{
		Ta có $A=\left\{(1,6);(2,6);(3,6);(4,6);(5,6);(6,6);(6,1);(6,2);(6,3);(6,4);(6,5)\right\}$.
	}
\end{ex}

%Câu 15
\begin{ex}%[0T8B2-2]%[Dự án đề kiểm tra HKII NH22-23- Phạm Văn Long]%[Phan Ngọc Hiển - Cà Mau]
	Có $8$ quả ổi và $6$ quả xoài. Có bao nhiêu cách chọn ra một quả trong các quả ấy?
	\choice
	{$24$}
	{$48$}
	{\True $14$}
	{$18$}
	\loigiai{
		Số cách chọn $1$ quả trong các quả trên là $\mathrm{C}^1_{14} = 14$ cách.
	}
\end{ex}

%Câu 16
\begin{ex}%[0T0B1-1]%[Dự án đề kiểm tra HKII NH22-23- Phạm Văn Long]%[Phan Ngọc Hiển - Cà Mau]
	Gieo $3$ đồng tiền là một phép thử ngẫu nhiên có không gian mẫu là
	\choice
	{$\left\{NNN; SSS; NNS; SSN; NSS; SNN\right\}$}
	{$\left\{NN; NS ; SN; SS\right\}$}
	{$\left\{NNN; SSS; NNS ;SSSN ;NSN; SNS\right\}$}
	{\True $\left\{NNN; SSS; NNS; SSN; NSN; SNS; NSS; SNN \right\}$}
	\loigiai{
		Ta có không gian mẫu của phép thử ngẫu nhiên trên là $\left\{NNN; SSS; NNS; SSN; NSN; SNS; NSS; SNN \right\}$.
	}
\end{ex}

%Câu 17
\begin{ex}%[0T9B2-1]%[Dự án đề kiểm tra HKII NH22-23- Phạm Văn Long]%[Phan Ngọc Hiển - Cà Mau]
	Đường thẳng $\Delta$ đi qua điểm $M(1;4)$ và có véc-tơ pháp tuyến $\overrightarrow{n}=(2;3)$ có phương trình tổng quát là
	\choice
	{$-x+4y+10=0$}
	{$2x+3y + 10=0$}
	{$-x+4y -10 =0$}
	{\True $2x +3y-14=0$}
	\loigiai{
		Phương trình tổng quát của đường thẳng $\Delta$ đi qua điểm $M(1;4)$ và có véc-tơ pháp tuyến $\overrightarrow{n}=(2;3)$ là
		\[2(x-1) + 3(y-4) = 0 \Leftrightarrow 2x +3y -14=0.\]
	}
\end{ex}

%Câu 18
\begin{ex}%[0T9B3-1]%[Dự án đề kiểm tra HKII NH22-23- Phạm Văn Long]%[Phan Ngọc Hiển - Cà Mau]
	Tọa độ tâm $I$ và bán kính $R$ của đường tròn $(C)\colon (x-1)^2 + (y+3)^2 = 16$ là
	\choice
	{$I\left(-1;3\right), R=4$}
	{$I\left(1;-3\right), R=16$}
	{$I\left(-1;3\right), R=16$}
	{\True $I\left(1;-3\right), R=4$}
	\loigiai{
		Tọa độ tâm $I$ và bán kính $R$ của đường tròn $(C)\colon (x-1)^2 + (y+3)^2 = 16$ là $I(1;-3)$, $R=4$.
	}
\end{ex}

%Câu 19
\begin{ex}%[0T9B2-2]%[Dự án đề kiểm tra HKII NH22-23- Phạm Văn Long]%[Phan Ngọc Hiển - Cà Mau]
	Cho $\triangle ABC$ có $A(2;-1)$, $B(4;5)$, $C(-3;2)$. Viết phương trình tổng quát của đường cao $AH$.
	\choice
	{\True $7x + 3y -11 = 0$}
	{$-3x + 7y +13 = 0$}
	{$7x + 3y +13 = 0 $}
	{$3x + 7y + 1=0$}
	\loigiai{
		Ta có $\overrightarrow{BC} = \left(-7;-3\right)$.\\
		Đường cao $AH$ qua $A(2;-1)$ và nhận $\overrightarrow{BC} = \left(-7;-3\right)$ làm véc-tơ pháp tuyến nên có phương trình tổng quát là
		\[(BC)\colon -7(x-2) -3(y+1) = 9 \Leftrightarrow -7x -3y +11 = 0 \Leftrightarrow 7x+3y-11=0.\]
	}
\end{ex}

%Câu 20
\begin{ex}%[0T9B3-2]%[Dự án đề kiểm tra HKII NH22-23- Phạm Văn Long]%[Phan Ngọc Hiển - Cà Mau]
	Cho hai điểm $A(5;-1)$, $B(-3;7)$. Đường tròn có đường kính $AB$ có phương trình là
	\choice
	{$(C)\colon (x+1)^2 + (y+3)^2 = 128$}
	{$(C)\colon (x+1)^2 + (y+3)^2 = 32$}
	{$(C)\colon (x-1)^2 + (y-3)^2 = 128$}
	{$(C)\colon (x-1)^2 + (y-3)^2 = 32$}
	\loigiai{
		Ta có $AB=\sqrt{(-3-5)^2 + (7+1)^2} = 8\sqrt{2}$.\\
		Gọi $I$ là trung điểm $AB$. Suy ra $I\left(1;3\right)$.\\
		Khi đó đường tròn đường kính $AB$ có tâm $I(1;3)$ và $R=\dfrac{AB}{2} = 4\sqrt{2}$ nên có phương trình là
		\[(C)\colon (x-1)^2 + (y-3)^2 = 32.\]
	}
\end{ex}
%Câu 21...........................
\begin{ex}%[0T9B2-2]%[Dự án đề kiểm tra HKII NH22-23- Phạm Văn Long]%[Phan Ngọc Hiển - Cà Mau]
Véc-tơ chỉ phương của đường thẳng $d\colon \heva{&x=1-4 t\\&y=-2+3 t}$ là
	\choice
	{\True $\vec{u}=(-4 ; 3)$}
	{$\vec{u}=(3 ; 4)$}
	{$\vec{u}=(1 ;-2)$}
	{$\vec{u}=(4 ; 3)$}
	\loigiai{
	Véc-tơ chỉ phương của đường thẳng $d$ là $\vec{u}=(-4 ; 3)$.	
	}
\end{ex}
%Câu 22...........................
\begin{ex}%[0T8Y2-1]%[Dự án đề kiểm tra HKII NH22-23- Phạm Văn Long]%[Phan Ngọc Hiển - Cà Mau]
	Khẳng định nào sau đây đúng?
	\choice
	{\True $\mathrm{C}_n^k=\dfrac{n !}{k !(n-k) !}$}
	{$\mathrm{C}_n^k=\dfrac{n !}{(n-k) !}$}
	{$\mathrm{C}_n^k=\dfrac{k !}{n !(n-k) !}$}
	{$\mathrm{C}_n^k=\dfrac{k !}{(n-k) !}$}
	\loigiai{
	Khẳng định nào sau đây đúng là $\mathrm{C}_n^k=\dfrac{n !}{k !(n-k) !}$.	
	}
\end{ex}
%Câu 23...........................
\begin{ex}%[0T8B3-2] %[Dự án đề kiểm tra HKII NH22-23- Phạm Văn Long]%[Phan Ngọc Hiển - Cà Mau]
	Trong khai triển nhị thức Niu-tơn của $(2x+5)^4$ hệ số của $x^3$ bằng
	\choice
	{$16$}
	{$600$}
	{$150$}
	{\True $160$}
	\loigiai{
	Ta có $(2x+5)^4=(2x)^4+4\cdot(2x)^3\cdot 5+6\cdot(2x)^2\cdot 5^2+4\cdot(2x)\cdot 5^3+5^4$.\\
	Suy ra hệ số của $x^3$ bằng $4\cdot 2^3\cdot 5=160$.
	}
\end{ex}
%Câu 24...........................
\begin{ex}%[0T7Y2-1]%[Dự án đề kiểm tra HKII NH22-23- Phạm Văn Long]%[Phan Ngọc Hiển - Cà Mau]
	Bất phương trình nào sau đây là bất phương trình bậc hai một ẩn
	\choice
	{$2 x+1>0$}
	{\True $x^2+3 x+6>0$}
	{$2 x+3 y<0$}
	{$3 x^2+4 y^2<0$}
	\loigiai{
	Bất phương trình $x^2+3 x+6>0$ là bất phương trình bậc hai một ẩn.	
	}
\end{ex}
%Câu 25...........................
\begin{ex}%[0T7B3-2]%[Dự án đề kiểm tra HKII NH22-23- Phạm Văn Long]%[Phan Ngọc Hiển - Cà Mau]
	Phương trình $\sqrt{x-1}=x-3$ có tập nghiệm là
	\choice
	{$S=\{2 ; 5\}$}
	{$S=\varnothing$}
	{$S=\{2\}$}
	{\True $S=\{5\}$}
	\loigiai{
	Ta có $\sqrt{x-1}=x-3\Rightarrow x-1=\left(x-3\right)^2\Rightarrow x^2-7x+10=0\Rightarrow \hoac{&x=2\\&x=5.}$\\
	Thay vào phương trình ta thấy $x=5$ thỏa mãn.
	}
\end{ex}
%Câu 26...........................
\begin{ex}%[0T9B3-2] %[Dự án đề kiểm tra HKII NH22-23- Phạm Văn Long]%[Phan Ngọc Hiển - Cà Mau]
	Đường tròn có tâm $I(1 ; 2)$, bán kính $R=3$ có phương trình là
	\choice
	{$(C)\colon (x+1)^2+(y+2)^2=9$}
	{$(C)\colon (x+1)^2+(y+2)^2=3$}
	{\True $(C)\colon (x-1)^2+(y-2)^2=9$}
	{$(C)\colon (x-1)^2+(y-2)^2=3$}
	\loigiai{
	Đường tròn có tâm $I(1 ; 2)$, bán kính $R=3$ có phương trình là	$(C)\colon (x-1)^2+(y-2)^2=9$.
	}
\end{ex}
%Câu 27...........................
\begin{ex}%[0T9Y4-2]%[Dự án đề kiểm tra HKII NH22-23- Phạm Văn Long]%[Phan Ngọc Hiển - Cà Mau]
	Phương trình nào sau đây là phương trình chính tắc của một elip?
	\choice
	{\True $\dfrac{x^2}{5}+\dfrac{y^2}{2}=1$}
	{$\dfrac{x^2}{4}+\dfrac{y^2}{25}=-1$}
	{$\dfrac{x^2}{4}+\dfrac{y^2}{25}=0$}
	{$\dfrac{x^2}{4}-\dfrac{y^2}{25}=1$}
	\loigiai{
	Phương trình $\dfrac{x^2}{5}+\dfrac{y^2}{2}=1$ là phương trình chính tắc của một elip.	
	}
\end{ex}
%Câu 28...........................
\begin{ex}%[0T7B2-1] %[Dự án đề kiểm tra HKII NH22-23- Phạm Văn Long]%[Phan Ngọc Hiển - Cà Mau]
	Tập nghiệm của bất phương trình $x^2-4 x-5 \leq 0$ là
	\choice
	{$S=(-\infty ;-1) \cup(5 ;+\infty)$}
	{$S=(-\infty ;-1] \cup[5 ;+\infty)$}
	{\True $S=[-1 ; 5]$}
	{$S=(-1 ; 5)$}
	\loigiai{
	Ta có $x^2-4 x-5 \leq 0\Leftrightarrow -1\le x \le 5$.\\
	Vậy $S=[-1 ; 5]$.	
	}
\end{ex}


\Closesolutionfile{ans}
%\begin{center}
%	\textbf{ĐÁP ÁN}
%	\inputansbox{10}{ans/ans}	
%\end{center}
\begin{center}
	\textbf{PHẦN 2 - TỰ LUẬN}
\end{center}
\begin{bt}%[0T7B2-1]%[Dự án đề kiểm tra HKII NH22-23- Phạm Văn Long]%[Phan Ngọc Hiển - Cà Mau]
	Giải bất phương trình $x^2-3 x+2>0$.
\loigiai{
Cho $x^2-3 x+2=0 \Leftrightarrow\hoac{&x=1\\&x=2.}$\\
Lập bảng xét dấu
	\begin{center}
	\begin{tikzpicture}
		\tkzTabInit[lgt=2.5,espcl=2.5,deltacl=0.6]
		{$x$ /0.6,$x^2-3 x+2$ /0.6}
		{$-\infty$,$1$,$2$,$+\infty$}
		\tkzTabLine{,+,$0$,-,$0$,+,}		
	\end{tikzpicture}
\end{center}
Suy ra tập nghiệm $T=(-\infty ; 1) \cup(2 ;+\infty)$.
}
\end{bt}
\begin{bt}%[0T0B2-3]%[Dự án đề kiểm tra HKII NH22-23- Phạm Văn Long]%[Phan Ngọc Hiển - Cà Mau]
	Một hộp chứa 5 quả cầu màu xanh và 6 quả cầu màu đỏ. Chọn ngẫu nhiên $2$ quả cầu. Tính xác suất để $2$ quả cầu được chọn cùng màu.
	\loigiai{
Số phần tử không gian mẫu: $n(\Omega)=C_{11}^2=55$.\\
Gọi $\mathrm{A}$ là biến cố" Hai quả cầu được chọn cùng màu".
$$n(A)=C_5^2+C_6^2=25.$$
Xác suất cần tìm là $p(A)=\dfrac{25}{55}=\dfrac{5}{11}$.
}
\end{bt}
\begin{bt}%[0T9B3-2]%[Dự án đề kiểm tra HKII NH22-23- Phạm Văn Long]%[Phan Ngọc Hiển - Cà Mau]
	Trong mặt phẳng tọa độ $Oxy$, cho điểm $A(1;3)$ và $B(2;5)$.
	\begin{enumerate}
		\item Viết phương trình đường thẳng $d$ đi qua $A$ và $B$.
		\item Viết phương trình đường tròn $(C)$ có tâm $A$ đi qua $B$.
	\end{enumerate}
	\loigiai{
\begin{enumerate}
	\item Ta có $\overrightarrow{AB}=(1;2)$.\\
	Phương trình tham số đường thẳng $d$ đi qua $A$ và $B$ là 
	$\heva{&x=1+t\\&y=3+2t.}$
	\item Ta có $\overrightarrow{AB}=(1;2) \Rightarrow|\overrightarrow{AB}|=\sqrt{5}$.\\
	Suy ra phương trình đường tròn $(C)$ có tâm $A$ đi qua $B$ là $(C)\colon (x-1)^2+(y-3)^2=5$.
\end{enumerate}
}
\end{bt}
\begin{bt}% [0T8K2-5]%[Dự án đề kiểm tra HKII NH22-23- Phạm Văn Long]%[Phan Ngọc Hiển - Cà Mau]
Cần sắp xếp thứ tự 8 tiết mục văn nghệ cho buổi biểu diễn văn nghệ của trưởng Ban tổ chức dự kiến xếp $4$ tiết mục ca nhạc ở vị trí thứ $1$, thứ $2$, thứ $5$ và thứ $8$; $2$ tiết mục múa ở vị trí thứ $3$ và thứ $6$; $2$ tiết mục hài ờ vị trí thứ $4$ và thứ $7$ . Có bao nhiêu cách xếp khác nhau?
	\loigiai{
		\begin{itemize}
			\item Sắp xếp $4$ tiết mục ca nhạc vào $4$ vị trí ($1,2,5$ và $8$) có $4!$ cách.
			\item Sắp xếp $2$ tiết mục múa vào $2$ vị trí ($3$ và $6$) có $2!$ cách.
			\item Sắp xếp $2$ tiết mục hài vào $2$ vị trí ($4$ và $7$) có $2!$ cách.
		\end{itemize}
		Vây có $4!\cdot 2!\cdot 2!=96$ cách.
}
\end{bt}